\section{Borsuk--Ulam theorem}

\begin{definition}[Sphere and antipodal map]
\label{def:bu-antipode}
\lean{BorsukUlamTheorem.UnitSphere,BorsukUlamTheorem.antipode}
\leanok
Let $S^n\subset\mathbb R^{n+1}$ be the unit sphere and let $a(x)=-x$ be its antipodal map.
\end{definition}

\begin{lemma}[Antipodal map properties]
\label{lem:bu-antipode}
\lean{BorsukUlamTheorem.continuous_antipode,BorsukUlamTheorem.antipode_antipode}
\leanok
\uses{def:bu-antipode}
The antipodal map is continuous and involutive.
\end{lemma}

\begin{lemma}[Odd maps on connected spheres]
\label{lem:bu-odd-zero}
\lean{BorsukUlamTheorem.odd_map_zero_of_connected}
\leanok
\uses{lem:bu-antipode}
A continuous odd real-valued function on a connected positive-dimensional sphere has a zero.
\end{lemma}

\begin{theorem}[Borsuk--Ulam]
\label{thm:borsuk-ulam}
\lean{BorsukUlamTheorem.MainTheorem}
\notready
\uses{lem:bu-antipode,lem:bu-odd-zero}
Every continuous map $S^n\to\mathbb R^n$ identifies a pair of antipodal points. The Lean file proves the low-dimensional cases; the case $n\ge2$ still contains a \texttt{sorry}.
\end{theorem}
